% --- Archon physics formalization source begin ---
% archon:physics
% archon:chemistry
% archon:covers IChO2026Problems/problem_icho_2026_t4_a1.lean
% archon:source-report reports/icho_2026/problem_icho_2026_t4_a1.source.json
% archon:problem-id icho_2026_t4
% archon:part-id T4-A1

\chapter{Icho Chemistry problem icho\_2026\_t4\_a1}
\label{ch:IChO2026Problems_problem_icho_2026_t4_a1}

\paragraph{Problem source.}
\#\# Source
58th International Chemistry Olympiad, Tashkent, Uzbekistan, 2026.
Official-materials index: https://www.icho2026.uz/problems
English problem PDF: ../icho\_2026\_source/raw/theory\_problem.pdf
English solution/rubric PDF: ../icho\_2026\_source/raw/theory\_solution.pdf
Problem PDF page 37; printed local question page 1; solution starts on PDF page 32.
The page PNG is primary visual evidence for formulas, structures, tables, and figures.

\#\# T4. The Nuclear Past of Uzbekistan
T4. The nuclear past of Uzbekistan 6\% Natural uranium consists primarily of two isotopes, 235U and 238U . 235U is responsible for sustaining nuclear fission reactions.

\#\# Current subquestion 4.1
Calculate the atomic abundance of 235U in natural uranium, assuming it consists only of isotopes 235U (235.04 a.u.) and 238U (238.05 a.u.).

\#\# Dataset metadata
IChO 2026; T4; raw rubric points=2; kind=theory; formalization\_ready=true.

\paragraph{Shared problem context.}
T4. The nuclear past of Uzbekistan 6\% Natural uranium consists primarily of two isotopes, 235U and 238U . 235U is responsible for sustaining nuclear fission reactions.

\paragraph{Current subquestion.}
Calculate the atomic abundance of 235U in natural uranium, assuming it consists only of isotopes 235U (235.04 a.u.) and 238U (238.05 a.u.).

\paragraph{Recorded answer/context.}
Using the average atomic mass of uranium from the Periodic Table as 238.03 a.u., and isotope masses as 235.04 a.u. for 235U and 238.05 a.u. for 238U, we can calculate the atomic abundance, 𝑥, of 235U. Equation: 238.03 = 𝑥× 235.04 + (1 −𝑥) × 238.05 (1) 238.03 = 235.04𝑥+ 238.05 −238.05𝑥 238.03 −238.05 = 𝑥(235.04 −238.05) −0.02 = 𝑥(−3.01) 𝑥= −0.02 −3.01 = 0.00664 = 0.66\% The atomic abundance of 235U in natural uranium is approximately 0.66\%, which is a slightly lower estimate compared to the commonly accepted value of about 0.72\%, influenced by rounding and exact values used. The atomic abundance of 235U is 0.66\% 2 points total: 1 point for equation (1); 1 point for the correct answer. 2.0 pt

\paragraph{Figure/image path.}
../icho\_2026\_source/image/T4\_page-1.png

\paragraph{Formalization target.}
The verified Lean formalization is in `IChO2026Problems/problem\_icho\_2026\_t4\_a1.lean`,
with its theorem and lemma proofs complete.
The Lean declarations must preserve every mathematical object, quantifier, hypothesis, side condition, approximation/error guarantee, and requested conclusion from the source statement.
The implementation reuses the pinned Mathlib, Physlib, and Chemistry libraries,
with local definitions only for source-specific objects.

\begin{definition}[Uranium isotopes, mass readouts, and isotopic composition]
\label{def:physics:icho_2026_t4_a1:setup}
\lean{IChO2026Problems.T4.A1.UraniumIsotope, IChO2026Problems.T4.A1.UraniumIsotope.mass, IChO2026Problems.T4.A1.IsotopicComposition, IChO2026Problems.T4.A1.naturalUraniumAverageMass, IChO2026Problems.T4.A1.averageAtomicMass}
\leanok
Natural uranium is modeled as a sample consisting only of the two isotopes
${}^{235}\mathrm{U}$ and ${}^{238}\mathrm{U}$, with printed isotope masses
$235.04$ a.u.\ and $238.05$ a.u.  An isotopic composition assigns each
isotope its atom fraction, the periodic-table average atomic mass of natural
uranium is $238.03$ a.u., and the average-mass (mixing) law takes the
atom-fraction-weighted sum of the isotope masses.
\end{definition}

\begin{theorem}[Atomic abundance of ${}^{235}\mathrm{U}$ in natural uranium]
\label{thm:physics:icho_2026_t4_a1:target}
\lean{IChO2026Problems.T4.A1.atomic_abundance_u235}
\leanok
If an isotopic composition of natural uranium assigns a nonnegative atom
fraction to each of the two isotopes, is normalized (the sample consists only
of ${}^{235}\mathrm{U}$ and ${}^{238}\mathrm{U}$), and reproduces the
periodic-table average atomic mass $238.03$ a.u.\ under the average-mass law,
then the atomic abundance of ${}^{235}\mathrm{U}$ is exactly
$0.02 / 3.01 \approx 0.00664$, i.e.\ $0.66\,\%$ at the precision quoted by
the official answer ($|100 x - 0.66| \le 0.005$).
\end{theorem}
\begin{proof}
\leanok
Write $x$ for the ${}^{235}\mathrm{U}$ atom fraction.  Normalization gives
the ${}^{238}\mathrm{U}$ fraction $1 - x$, and substituting the printed
isotope masses and the periodic-table average mass into the mixing law yields
equation (1) of the rubric,
$238.03 = 235.04\,x + 238.05\,(1 - x)$.
Collecting the linear terms,
$x\,(235.04 - 238.05) = 238.03 - 238.05$, so $-3.01\,x = -0.02$ and
$x = 0.02 / 3.01 \approx 0.00664$.
As a percentage, $100 x = 200 / 301 \approx 0.6645$, which differs from the
official $0.66\,\%$ readout by less than $0.005$.
In Lean the equation is solved by \texttt{linarith} after unfolding the
average-mass law, and the percentage bound is closed by \texttt{norm_num}
via \texttt{abs\_le}.
\end{proof}
% --- Archon physics formalization source end ---
